%------------------------------------------------------------------------
% 书脊版式【学生不需要改动】
%
% 按《广东工业大学毕业设计（论文）手册》2025 版样张 1(2)：
%   - 整页居中灰色书脊条
%   - 论文题目：四号黑体加粗，距页面顶端 20 mm
%   - 学生姓名：四号黑体加粗，位于题目与学院名之间
%   - 学院名称：小四号黑体加粗，距页面底端 20 mm
%   - 所有字符竖排（逐字换行，紧贴排列）
%
% 题目、姓名、学院名均取自根目录的 学生信息.tex，不用在本文件里改。
% 书脊通常由打印社单独裁切，这里以整页 A4 输出供参考。
%
% 说明：
%   - 使用临时零页边距让灰色书脊条覆盖整页高度。
%   - 页面输出后恢复原页边距，但不额外换页，避免书脊后插入空白页。
%   - 三个 minipage 内部强制 baselineskip = 字号，避免继承正文 1.5 倍行距。
%   - 两段 \vfill 负责均匀分配题目、姓名、学院名之间的空白。
%------------------------------------------------------------------------
\newgeometry{margin=0pt}
\thispagestyle{empty}

\newlength{\gdutspinewidth}
\setlength{\gdutspinewidth}{13mm}

\begingroup
\setlength{\parindent}{0pt}
\setlength{\fboxsep}{0pt}
\makeatletter

\noindent\makebox[\paperwidth]{%
  \colorbox{black!20}{%
    \parbox[t][\paperheight][t]{\gdutspinewidth}{%
      \centering
      % —— 顶部留白 20 mm ——
      \vspace*{20mm}

      % 论文题目（四号黑体加粗，baseline = 1em，逐字竖排）
      {\fontsize{14bp}{14bp}\heiti\bfseries\selectfont
       \begin{minipage}[t]{1em}\centering\@title\end{minipage}\par}

      \vfill

      % 学生姓名（四号黑体加粗，baseline = 1em，逐字竖排）
      {\fontsize{14bp}{14bp}\heiti\bfseries\selectfont
       \begin{minipage}[t]{1em}\centering\@author\end{minipage}\par}

      \vfill

      % 学院名称（小四号黑体加粗，baseline = 1em，逐字竖排）
      {\fontsize{12bp}{12bp}\heiti\bfseries\selectfont
       \begin{minipage}[t]{1em}\centering\@college\end{minipage}\par}

      % —— 底部留白 20 mm ——
      \vspace*{20mm}
    }%
  }%
}\par

\makeatother
\endgroup
\clearpage

\makeatletter
\Gm@restore@pkg
\Gm@changelayout

% 书脊页已经通过上面的 \clearpage 结束；下一页通常从 \chapter* 摘要开始。
% 跳过紧随其后的那一次 \clearpage，避免在书脊和摘要之间产生空白页。
\let\gdutspineorigclearpage\clearpage
\gdef\clearpage{%
  \global\let\clearpage\gdutspineorigclearpage}
\makeatother
